\chapter{Literature Review}
\label{ch:literature}

\section{WiFi Sensing and Channel State Information}

In an OFDM system the transmitted signal is divided across many narrowband
subcarriers. The receiver estimates, for each subcarrier $k$, a complex
channel response $H_k = |H_k| e^{j\angle H_k}$ relating the transmitted and
received symbols. The collection of these estimates across subcarriers is the
Channel State Information.

The physical basis for sensing is multipath. The received signal is the sum
of components arriving along paths of differing length; a path of length $d$
contributes a phase rotation proportional to $d/\lambda$, so at 2.4\,GHz,
where $\lambda \approx 12.5$\,cm, a path-length change of a few centimetres
produces a full phase rotation and a substantial amplitude change through
constructive and destructive interference. A moving human body changes many
path lengths simultaneously, and the resulting time-varying pattern across
subcarriers carries a signature of the motion. Because path-length change
maps to phase change differently at each subcarrier frequency, the pattern is
frequency-selective --- which is precisely why CSI supports finer
discrimination than the single RSSI scalar, which collapses all of this into
one number.

Miao et al.\ \cite{miao2025survey} survey the field and organise it into a
common pipeline --- acquisition, denoising, feature extraction,
classification --- that this project follows. Their survey also documents the
field's persistent difficulty: models generalise poorly across environments
and subjects, because the learned mapping depends on a specific geometry of
reflectors rather than on any subject-independent property of the activity.

\section{Deep Learning for CSI-Based HAR}

Early CSI recognition systems used hand-crafted statistical and spectral
features with classical classifiers. Deep models displaced them by learning
representations directly from the CSI time--frequency matrix. Convolutional
networks treat the (time $\times$ subcarrier) matrix as an image and capture
localised patterns; recurrent architectures, particularly LSTMs, model the
temporal evolution of an activity; hybrid CNN-LSTM designs use a
convolutional encoder for spectral structure and a recurrent layer for
temporal dependency. More recently transformer encoders have been applied to
the same tensors.

The SenseFi benchmark of Yang et al.\ \cite{yang2023sensefi} is the most
useful reference point, providing a PyTorch library and a common evaluation
of these model families across public WiFi sensing datasets. Two of its
findings inform the design here. Model architecture matters less than is
often implied --- differences between well-tuned families are frequently
smaller than differences between datasets or preprocessing choices --- and
reported numbers depend heavily on the evaluation regime. This project
accordingly limits itself to a CNN and a CNN-LSTM, spending its effort on
data quality and evaluation rather than on architecture search.

A methodological caution runs through this literature and directly motivates
this project's protocol. Many reported results, including strong LSTM figures
obtained by $k$-fold cross-validation, assign \emph{windows} randomly to
folds. When windows overlap, or even when they merely come from the same
continuous recording, temporally adjacent and nearly identical windows land
on both sides of the split. The resulting accuracy measures how well a model
recognises a recording, not an activity, and it is an unreliable predictor of
deployed behaviour. Section~\ref{sec:eval-protocol} describes the three-split
protocol adopted in response.

\section{HAR on ESP32-Class Hardware}

Access to CSI on inexpensive microcontrollers changed what is practical
outside a laboratory. Hernandez and Bulut's ESP32-CSI-Tool was the first
widely used firmware for this purpose and established the basic approach of
streaming CSI over serial to a host; it targets ESP-IDF v4.3 and the original
ESP32, and remains the reference cited by much of the ESP32 sensing
literature. Espressif's own \texttt{esp-csi} framework \cite{espcsi} now
provides maintained support across the chip family, including the ESP32-S3,
along with an example matching a router-plus-receivers topology. Espressif
additionally ranks CSI quality across its chips, placing the S3 in the middle
of the range and above the original ESP32 used in most published ESP32
activity-recognition work.

Moshiri et al.\ \cite{moshiri2021csihar} released a public ESP32 CSI dataset
covering seven activities and demonstrated that deep models achieve high
accuracy on this class of hardware, establishing the basic feasibility this
project builds on. Their dataset uses a single receiver. Strohmayer and
Kampel \cite{strohmayer2023wall} showed long-range through-wall recognition
with ESP32 hardware, confirming that the constraint is not the radio's
sensitivity. The more recent ESP-Fi dataset \cite{espfi2026} continues this
direction.

Two consistent properties of this body of work shaped the present design.
Nearly all of it is amplitude-only, for the single-antenna reason set out in
Chapter~\ref{ch:introduction}. And most of it uses a single receiver, so the
benefit of multiple spatially separated receivers --- a central claim of this
project --- is asserted more often than it is measured, which is why the
receiver-count ablation is reported here as a primary result rather than as
an afterthought.

\section{Research Gap}

Three gaps follow from the above.

\textbf{Honest generalisation measurement on low-cost hardware.} Feasibility
on ESP32-class devices is established, but reported figures are dominated by
random-split evaluation. What a deployment needs to know --- how much
accuracy is lost on an unseen day, and on an unseen person --- is rarely
reported for this hardware class. This project reports all three regimes for
every model, and treats the gap between them as a result in its own right
rather than as a number to minimise.

\textbf{Multi-receiver benefit, quantified.} Multiple receivers are commonly
justified on the intuition that complementary geometry helps. This project
measures the contribution directly by ablating receiver count on identical
data with an architecture that admits the change without modification.

\textbf{Falls as a first-class metric.} Fall detection is the usual
motivation for this work, yet falling is typically reported inside an
aggregate accuracy that its scarcity allows it to barely influence. This
project reports falling recall separately in every result, and treats a
system that scores well overall while missing falls as having failed at its
stated purpose.

A locally collected multi-receiver ESP32-S3 dataset, with a datasheet and
per-session metadata, is a secondary contribution: no comparable multi-receiver
dataset for this hardware is publicly available, and the collection
context --- a Nepali deployment setting with commodity equipment --- differs
from that of the existing public datasets.
